% Standalone contribution; requires amsmath, amssymb, amsthm, booktabs, hyperref.
% Bibliography entries are supplied in depth_cayley_references.tex.
\providecommand{\shuffle}{\mathbin{\sqcup\!\sqcup}}
\section{A Cayley normal form that preserves series depth}
\label{depthcay:sec}

The inspected continuation already supplies a constructive projector for
 the full shuffle ideal generated by convergent Cayley identities
\cite{depth:PriorProjection}. Its image is fixed by the Cayley involution,
but its expansion can increase the number of nonzero letters. We construct
another projector for the same ideal whose output never has greater
series depth than its input. The choice of an oriented representative,
and a different equivariant endpoint projection, make this possible.
We then compute the entire depth filtration of that formal quotient.
These are statements about a specified algebra of words. The period map
may satisfy further identities, so no numerical independence assertion is
involved.

The underlying shuffle-polynomial theorem, Eulerian logarithm, and Witt
character formula are classical \cite{depth:Radford,depth:Patras,depth:Reutenauer}.
The contribution here is their simultaneous compatibility with the two
letter counts exchanged by the Cayley transformation, the resulting
non-increasing-depth retraction, and the explicit filtered dimensions.

\subsection{Alphabet, admissibility, and the two counts}

Use outer-letter-first iterated integrals on $[0,1]$ and the forms
\[
 z=\frac{dt}{t},\qquad c=\frac{dt}{1-t},\qquad
 a=\frac{i\,dt}{1-it},\qquad b=-\frac{dt}{1+t},\qquad
 \bar a=-\frac{i\,dt}{1+it}.
\]
An ordinary convergent word has first letter different from $c$ and last
letter different from $z$. Introduce the adapted letters
\begin{equation}\label{depthcay:eq:alphabet}
 x=c-b,\qquad y=a-\bar a,\qquad h=a+\bar a-b,
 \qquad \mathcal E=\{z,y,h,b,x\}.
\end{equation}
They form a rational basis of the same five-dimensional space of forms.
In particular,
\[
 c=x+b,\qquad a=\frac{y+h+b}{2},\qquad
 \bar a=\frac{-y+h+b}{2}.
\]
Every adapted letter other than $z$ is a linear combination of nonzero
original letters. Hence the original series-depth filtration is exactly
\begin{equation}\label{depthcay:eq:depth}
 \operatorname{dep}(w)=|w|-n_z(w),\qquad
 F_D\mathcal A=\operatorname{span}_{\mathbb Q}
 \{w:\operatorname{dep}(w)\le D\},
\end{equation}
where $n_z(w)$ counts occurrences of $z$.
Let $\widetilde{\mathcal A}$ be the full shuffle algebra on $\mathcal E$,
with empty word $1$ and product $\shuffle$. Let $\mathcal A$ be its
subalgebra spanned by $1$ and the adapted words satisfying
\begin{equation}\label{depthcay:eq:admissible}
 w_1\ne x,\qquad w_{|w|}\ne z.
\end{equation}
This is precisely the original convergent word space: $x$ is the only
adapted letter with a $c$ component, and the other four adapted letters
span the original first-letter subspace that excludes $c$.

The Cayley operator and complex conjugation are
\begin{align}
 C(v_1\cdots v_n)&=\tau(v_n)\cdots\tau(v_1),&
 \tau z=x,\quad\tau x=z,\quad\tau y=y,\quad
 \tau h=-h,\quad\tau b=b,\label{depthcay:eq:C}\\
 K(w)&=(-1)^{n_y(w)}w.&&\label{depthcay:eq:K}
\end{align}
They are commuting involutive shuffle-algebra automorphisms preserving
$\mathcal A$. The substitution $t\mapsto(1-t)/(1+t)$ has
$\tau=-\phi^*$; its orientation reversal is already included in
\eqref{depthcay:eq:C}. Thus, for the convergent period map $H$,
\begin{equation}\label{depthcay:eq:period-C}
 H(Cf)=H(f),\qquad H(f\shuffle g)=H(f)H(g),\qquad f,g\in\mathcal A.
\end{equation}
For example, $H(y)=i\pi/2$, $H(b)=-\log2$, and $H(h)=0$.
The first and last exclusions in \eqref{depthcay:eq:admissible} are
exchanged by reversal and $z\leftrightarrow x$, so every transformed
word remains convergent.

We impose exactly the full Cayley shuffle ideal
\begin{equation}\label{depthcay:eq:ideal}
 \mathcal J=\langle f-Cf:f\in\mathcal A\rangle_{\shuffle},
 \qquad \mathcal B=\mathcal A/\mathcal J.
\end{equation}
No stuffle, distribution, or independent period relation is included in
this definition. Equation \eqref{depthcay:eq:period-C} proves that the
period map annihilates $\mathcal J$.

For an adapted word use the bidegree
\[
 \deg_{z,x}(w)=(n_z(w),n_x(w)).
\]
Shuffle multiplication adds this bidegree. The operator $C$ exchanges
its two components, and $K$ preserves both.

\subsection{A count-preserving endpoint retraction}

\begin{lemma}[Homogeneous endpoint retraction]\label{depthcay:lem:R}
There is a unique shuffle-algebra retraction
$R_0:\widetilde{\mathcal A}\to\mathcal A$ fixing $\mathcal A$ and
sending the shuffle-polynomial variables $z,x$ to zero. It commutes
with $C,K$ and preserves the bidegree $(n_z,n_x)$ in every nonzero term.
If a word $w$ has $l$ initial $x$'s and $r$ terminal $z$'s, then
\begin{equation}\label{depthcay:eq:R}
 R_0(w)=\sum_{u=0}^{l}\sum_{v=0}^{r}(-1)^{u+v}
 x^u\shuffle w[u:|w|-v]\shuffle z^v.
\end{equation}
The powers and the middle subword in this formula use concatenation.
\end{lemma}

\begin{proof}
Order the adapted alphabet as $z<y<h<b<x$. Every Lyndon word of
length at least two is admissible. Indeed it cannot start with the
largest letter $x$, and cannot end with the smallest letter $z$.
The triangular Lyndon factorization and the shuffle-polynomial theorem
therefore give
\[
 \widetilde{\mathcal A}=\mathcal A[z,x].
\]
All these generators are homogeneous in the two counts. Setting $z,x$
to zero gives the unique retraction and preserves that homogeneity.
The ideal $(z,x)$ is stable under $C$ because $Cz=x$ and $Cx=z$,
and is stable under $K$. Both operators also preserve $\mathcal A$.
Their commutation with the retraction follows from this direct-sum
polynomial decomposition.

For the explicit formula, deleting one initial $x$ and deleting one
terminal $z$ are commuting derivations for shuffle multiplication.
Their Taylor substitutions set the corresponding polynomial variables
to zero. The $u$th shuffle power of one letter is $u!$ times its
concatenation power, canceling the Taylor factorial. This proves
\eqref{depthcay:eq:R}. Each of its terms before cancellation already
has the same two letter counts as $w$; after cancellation all remaining
words are admissible by the retraction property.
\end{proof}

This retraction is a formal polynomial operation. It assigns no analytic
value to divergent endpoint integrals. In the original coordinates its
substitution is $z=0$, $c=b$ in shuffle-polynomial coordinates, not a
letter-by-letter replacement inside a concatenation word. It is a
different equivariant retraction from the earlier choice
$z=-b/2$, $c=b/2$ \cite{depth:PriorProjection}. The present choice retains
the two homogeneous letter counts needed below.

\subsection{Eulerian lifts and the oriented projector}

For a nonempty word set
\begin{equation}\label{depthcay:eq:e}
 e(w)=\sum_{k=1}^{|w|}\frac{(-1)^{k-1}}{k}
 \sum_{\substack{w=w_1\cdots w_k\\|w_j|>0}}
 w_1\shuffle\cdots\shuffle w_k,
 \qquad e(1)=0.
\end{equation}
This is the first Eulerian logarithm for shuffle multiplication and
word deconcatenation. It annihilates products of positive-degree
factors and satisfies
\begin{equation}\label{depthcay:eq:e-properties}
 e^2=e,\qquad e(f)-f\in
 (\widetilde{\mathcal A}_{>0})^2
 \quad(f\in\widetilde{\mathcal A}_{>0}).
\end{equation}
For completeness, the product-annihilation property follows by viewing
$e$ as the convolution logarithm of the identity character into the
commutative shuffle algebra: that logarithm is an infinitesimal
character. The second statement is immediate from the $k=1$ summand
of \eqref{depthcay:eq:e}; product annihilation then gives $e^2=e$.
These are also the standard Eulerian-idempotent properties
\cite{depth:Patras}. Reversal permutes the consecutive cuts and reverses
the order of their commuting shuffle factors. Hence $e$ commutes with
$C$, as well as with $K$, directly from its displayed formula.
Every term preserves $(n_z,n_x)$.

Put $E=R_0e$ and $U=E(\mathcal A)$. The retraction $R_0$ and the
Eulerian lift both preserve the two counts, so $U$ is bigraded.
On a homogeneous element $u\in U$ of bidegree $(p,q)$ define
\begin{equation}\label{depthcay:eq:orientation}
 Q_{\uparrow}u=
 \begin{cases}
 u,&p>q,\\
 Cu,&p<q,\\
 (u+Cu)/2,&p=q.
 \end{cases}
\end{equation}
Thus a pair of exchanged components is represented on the side with
more copies of $z$. Let $E_{\uparrow}(w)=Q_{\uparrow}E(w)$; it can be
computed from the counts in $w$ itself, since $E$ preserves them.

\begin{theorem}[A non-increasing-depth normal form]\label{depthcay:thm:projector}
Define $\Pi_{\uparrow}(1)=1$ and
\begin{equation}\label{depthcay:eq:projector}
 \boxed{\displaystyle
 \Pi_{\uparrow}(w)=\sum_{k=1}^{|w|}\frac1{k!}
 \sum_{\substack{w=w_1\cdots w_k\\|w_j|>0}}
 E_{\uparrow}(w_1)\shuffle\cdots\shuffle E_{\uparrow}(w_k).}
\end{equation}
Then $\Pi_{\uparrow}:\widetilde{\mathcal A}\to\mathcal A$ is a graded
shuffle-algebra homomorphism. On $\mathcal A$ it satisfies
\begin{align}
 \Pi_{\uparrow}^2&=\Pi_{\uparrow},&
 \Pi_{\uparrow}C&=\Pi_{\uparrow},&
 \Pi_{\uparrow}K&=K\Pi_{\uparrow},\label{depthcay:eq:properties}\\
 \ker(\Pi_{\uparrow}|_{\mathcal A})&=\mathcal J,&
 \Pi_{\uparrow}(F_D\mathcal A)&\subseteq F_D\mathcal A.
 \label{depthcay:eq:kernel-depth}
\end{align}
Consequently
\begin{equation}\label{depthcay:eq:period-normalform}
 H(f)=H(\Pi_{\uparrow}f),\qquad f\in\mathcal A,
\end{equation}
and $f-g\in\mathcal J$ if and only if
$\Pi_{\uparrow}f=\Pi_{\uparrow}g$.
In general $C\Pi_{\uparrow}\ne\Pi_{\uparrow}$.
\end{theorem}

\begin{proof}
On $\mathcal A$, the map $E$ annihilates positive-degree products and
is congruent to the identity modulo such products. It is therefore
idempotent and its image $U$ maps isomorphically onto
$\mathcal A_{>0}/(\mathcal A_{>0})^2$. Homogeneous lifts of an
indecomposable basis freely generate the polynomial algebra
$\mathcal A$: the change from its Lyndon generators is triangular in
weight, with invertible linear part. Thus
\begin{equation}\label{depthcay:eq:SymU}
 \mathcal A\simeq\operatorname{Sym}(U)
\end{equation}
as a weight-graded and $(n_z,n_x)$-bigraded algebra. Both $C,K$ act
on $U$ because they commute with $R_0,e$.

For $p>q$, choose a $K$-eigenbasis of $U_{n,p,q}$ and the transported
basis $Cu$ of $U_{n,q,p}$. On $p=q$, diagonalize the commuting rational
involutions $C,K$. Equation \eqref{depthcay:eq:orientation} identifies
each transported pair, fixes the $C$-positive diagonal generators,
and kills the $C$-negative diagonal generators. It therefore extends
uniquely to an algebra retraction $P_{\uparrow}$ of
$\operatorname{Sym}(U)$. Its kernel is the ideal generated by
$u-Cu$. This ideal equals $\mathcal J$: one inclusion follows by
using $u\in\mathcal A$ in \eqref{depthcay:eq:ideal}; for the other,
$C$ acts trivially on the quotient in which these generator relations
are imposed. This proves the kernel statement and
$P_{\uparrow}C=P_{\uparrow}$. The selected eigenspaces also give
$P_{\uparrow}K=KP_{\uparrow}$.

The formula \eqref{depthcay:eq:projector} is exactly
$P_{\uparrow}R_0$. Indeed convolution exponentiation of
\eqref{depthcay:eq:e} gives, for every nonempty word,
\[
 w=\sum_{k=1}^{|w|}\frac1{k!}
 \sum_{w=w_1\cdots w_k}e(w_1)\shuffle\cdots\shuffle e(w_k).
\]
Apply $R_0$, then $P_{\uparrow}$. Each $E(w_j)$ belongs to $U$,
even if the block is not admissible. To justify this point, use
$\widetilde{\mathcal A}=\mathcal A[z,x]$: in degree at least two,
$w-R_0w$ is a sum of products of positive-degree polynomial factors
and is annihilated by $e$. In degree one the extra generators $z,x$
are killed by $R_0e$ directly. Thus $E(w)=E(R_0w)\in U$.
The resulting expression is exactly \eqref{depthcay:eq:projector},
proving multiplicativity and idempotence.

Finally each nonzero term of $E(w_j)$ has the original count pair
$(p_j,q_j)$. The orientation either retains $p_j$ copies of $z$,
replaces their number by $q_j>p_j$, or averages two terms with equal
counts. It never decreases the number of $z$'s. Shuffle multiplication
adds their counts, so every output word in
\eqref{depthcay:eq:projector} has at least $n_z(w)$ copies of $z$.
This proves the filtration statement word by word.

For admissible $f$, the kernel description gives
$f-\Pi_{\uparrow}f\in\mathcal J$, and
\eqref{depthcay:eq:period-C} proves the period identity.
To see the last assertion, $\Pi_{\uparrow}(zy)=zy$ but
$C\Pi_{\uparrow}(zy)=yx\ne zy$. Nevertheless
$\Pi_{\uparrow}(yx)=zy$, as required.
\end{proof}

One immediate all-weight example is
\[
 \Pi_{\uparrow}(yx^{n-1})=z^{n-1}y,\qquad
 H(yx^{n-1})=2i\,\beta(n),\qquad n\ge2.
\]
Here the second equality also follows from the single Cayley substitution
itself. The theorem deals with the full multiplicative ideal and applies
to arbitrary word polynomials, not just this example.

\subsection{The exact depth filtration of the full ideal quotient}

The filtration on $\mathcal B$ is the image of $F_D\mathcal A$.
It admits a completely explicit polynomial description.

\begin{theorem}[Filtered polynomial generators]\label{depthcay:thm:filtered}
Choose the $K$-eigenbases and $C$-paired bases in the proof of
Theorem~\ref{depthcay:thm:projector}. Then $\mathcal B$ is freely
polynomial on the following generators:
\begin{enumerate}
\item one representative of every $C$-paired basis vector in
$U_{n,p,q}\oplus U_{n,q,p}$, with $p>q$;
\item every $C$-positive basis vector of $U_{n,p,p}$.
\end{enumerate}
Their minimum depths in this formal quotient are respectively
$n-p=n-\max(p,q)$ and $n-p$. The depth of a monomial in these
polynomial generators is the sum of their listed depths. Consequently
these monomials give a basis adapted to every $F_D\mathcal B$.
\end{theorem}

\begin{proof}
Polynomial freeness follows from \eqref{depthcay:eq:SymU} by eliminating
one generator from each pair and every negative diagonal generator.
The polynomial coordinate changes are bigraded before the quotient.
For a surviving quotient monomial, one may independently lift each
paired factor on either side of its pair. Its greatest attainable
$z$ count is therefore the sum of the larger counts $\max(p,q)$.
The representative chosen above attains that count, and thus attains
the asserted depth.

Conversely, expand a homogeneous element of $F_D\mathcal A$ in the
bigraded polynomial basis of $\operatorname{Sym}(U)$. Every source
monomial has at least $n-D$ copies of $z$. A surviving quotient
monomial in its image has maximum attainable count no smaller than
that source count. Hence it belongs to the span of the basis monomials
whose listed depth is at most $D$. Quotient monomials are linearly
independent. A linear combination of lower-depth source elements
therefore cannot produce a basis monomial with larger listed depth.
This proves the assertion for every filtration step and the claimed
minimum within the formal quotient.
\end{proof}

We now give finite coefficient formulas for the number of these
polynomial generators. For $n\ge2$, $p,q\ge0$, $p+q\le n$, and
$\varepsilon\in\{1,-1\}$ put
\begin{equation}\label{depthcay:eq:W}
 W_{n,p,q}(\varepsilon)=\frac1n
 \sum_{d\mid\gcd(n,p,q)}\mu(d)
 \binom{n/d}{p/d}\binom{(n-p)/d}{q/d}
 (2+\varepsilon^d)^{(n-p-q)/d}.
\end{equation}
The convention is $\gcd(n,0,0)=n$.
Then the $K$-even and $K$-odd dimensions of $U_{n,p,q}$ are
\begin{equation}\label{depthcay:eq:a}
 a^+_{n,p,q}=\frac{W_{n,p,q}(1)+W_{n,p,q}(-1)}2,
 \qquad
 a^-_{n,p,q}=\frac{W_{n,p,q}(1)-W_{n,p,q}(-1)}2.
\end{equation}
For the diagonal Cayley traces set
\begin{equation}\label{depthcay:eq:T}
 T_n(\varepsilon;t)=\frac{(-1)^{n-1}}n
 \left[
 \sum_{\substack{d\mid n\\d\ {
m odd}}}\mu(d)\varepsilon^n
 +\sum_{\substack{d\mid n\\d\ {
m even}}}\mu(d)
       (2t^{d/2}+3)^{n/d}\right],
 \qquad T_{n,p}(\varepsilon)=[t^p]T_n(\varepsilon;t).
\end{equation}
Define
\[
 c^+_{n,p}=\frac{T_{n,p}(1)+T_{n,p}(-1)}2,\qquad
 c^-_{n,p}=\frac{T_{n,p}(1)-T_{n,p}(-1)}2.
\]
These are the traces of $C$ on the corresponding $K$-even and $K$-odd
parts of $U_{n,p,p}$. In particular,
\begin{equation}\label{depthcay:eq:oddtracezero}
 c^-_{n,p}=0\quad(n\ge2).
\end{equation}
Indeed for odd $n>1$ the entire polynomial \eqref{depthcay:eq:T}
vanishes by $\sum_{d\mid n}\mu(d)=0$; for even $n$ it is independent
of $\varepsilon$.

Let $g^\pm_{n,\delta}$ count the $K$-even/odd polynomial generators
of weight $n$ and minimum depth $\delta$. For $n\ge2$ they are
\begin{equation}\label{depthcay:eq:generators}
 g^\pm_{n,\delta}
 =\sum_{\substack{p>q\ge0\\p+q\le n\\n-p=\delta}}a^\pm_{n,p,q}
 +\sum_{\substack{2p\le n\\n-p=\delta}}
     \frac{a^\pm_{n,p,p}+c^\pm_{n,p}}2.
\end{equation}
At weight one the only surviving generators are $b$ and $y$, so
$g^+_{1,1}=g^-_{1,1}=1$.

\begin{proof}[Justification of the character formulas]
In weight $n\ge2$, deleting the degree-one shuffle generators $z,x$
does not change indecomposables. Their character is the free-Lie
character on the five-letter space. The classical Witt trace formula
\cite{depth:Reutenauer}, or equivalently PBW followed by formal logarithms
and M\"obius inversion, is
\[
 \operatorname{tr}(g\mid L_n(V))
 =\frac1n\sum_{d\mid n}\mu(d)
       \operatorname{tr}(g^d\mid V)^{n/d}.
\]
Apply it first to the diagonal operator with eigenvalues
$u,v,\varepsilon,1,1$. Extracting $u^pv^q$ gives
\eqref{depthcay:eq:W}; projection to the two $K$ eigenspaces gives
\eqref{depthcay:eq:a}.

On indecomposables, reversal acts by $(-1)^{n-1}$. This follows from
the shuffle antipode $S(w)=(-1)^n\operatorname{rev}(w)$ and the fact
that the antipode is minus the identity modulo positive-degree products.
Consequently $C$ has the character of $(-1)^{n-1}\tau$ there.
For the weighted operator combining $\tau$ with $K$ if
$\varepsilon=-1$, an odd power has trace $\varepsilon$ on the
one-letter space: the $z,x$ block has zero trace, and the $h,b$
contributions cancel. An even $d$th power has trace
$2(uv)^{d/2}+3$. Insert these expressions into the Witt trace formula
and put $t=uv$. Its coefficients are diagonal traces because a component
with $p\ne q$ is carried to a different bidegree. This is exactly
\eqref{depthcay:eq:T}. The same character calculation is valid on
the dual indecomposables: transpose has the same trace, including for
the weighted operator. Formula \eqref{depthcay:eq:generators} now
counts a full space for $p>q$ and the positive half, with its trace
correction, for $p=q$.
\end{proof}

The total Hilbert series and the $K$-trace series of the associated
depth-graded quotient are
\begin{align}
 \mathscr H(t,u)&=\prod_{n,\delta\ge1}
 (1-t^nu^\delta)^{-g^+_{n,\delta}-g^-_{n,\delta}},
 \label{depthcay:eq:Hilbert}\\
 \mathscr H_K(t,u)&=\prod_{n,\delta\ge1}
 (1-t^nu^\delta)^{-g^+_{n,\delta}}
 (1+t^nu^\delta)^{-g^-_{n,\delta}}.
 \label{depthcay:eq:HilbertK}
\end{align}
Thus the imaginary, or $K$-odd, dimensions are the coefficients of
$(\mathscr H-\mathscr H_K)/2$. All products are formal power series;
every fixed-weight coefficient is a finite calculation.
The first rows, by exact graded depth, are
\begin{center}
\begin{tabular}{c|rrrrrrr|r}
\toprule
weight & $1$ & $2$ & $3$ & $4$ & $5$ & $6$ & $7$ & total\\
\midrule
$1$ &1&&&&&&&1\\
$2$ &1&2&&&&&&3\\
$3$ &1&9&6&&&&&16\\
$4$ &1&16&40&15&&&&72\\
$5$ &1&22&116&161&42&&&342\\
$6$ &1&28&219&640&593&113&&1594\\
$7$ &1&34&347&1582&3115&2123&316&7518\\
\bottomrule
\end{tabular}
\end{center}
The total $7518$ agrees with the earlier full-ideal count
\cite{depth:PriorProjection}. Its refinement gives
\begin{equation}\label{depthcay:eq:weightseven}
 \dim F_4\mathcal B_7^-=1+34+347+1582=1964.
\end{equation}
The original imaginary admissible space of weight seven and depth at
most four has dimension $2546$ \cite{depth:PriorS6}; hence
\begin{equation}\label{depthcay:eq:intersection}
 \dim\bigl(\mathcal J_7^-\cap F_4\mathcal A_7\bigr)=582.
\end{equation}
These are dimensions of a formal quotient and of a formal ideal
intersection. They do not count linearly independent complex numbers.

\subsection{Completeness of the unmultiplied balanced rows}

One narrower consequence is useful for interpreting a bounded search.
Let $\mathcal L=(1-C)\mathcal A$ denote the linear span of unmultiplied
Cayley rows, without taking their ideal. For fixed weight $n$ and
$k=n-D\ge0$, let $B_{n,k}$ be the span of admissible adapted words
with at least $k$ copies of each of $z,x$.

\begin{proposition}\label{depthcay:prop:balanced}
\[
 \mathcal L_n\cap F_D\mathcal A_n=(1-C)B_{n,k}.
\]
Thus the balanced seeds exhaust all linear combinations of unmultiplied
Cayley rows whose resulting vector has depth at most $D$, even if
higher-depth terms were canceled before taking that result.
\end{proposition}
\begin{proof}
Since $C$ is an involution in characteristic zero,
$\mathcal L=\ker(1+C)$. A vector $r$ in the left side satisfies
$Cr=-r$, so both $r$ and $Cr$ belong to $F_D\mathcal A_n$.
The intersection $F_D\mathcal A_n\cap C(F_D\mathcal A_n)$ has
exactly the balanced adapted-word basis $B_{n,k}$. Therefore
$r=(1-C)(r/2)$ is in the right side. The reverse inclusion follows
because both a balanced word and its transform have depth at most $D$.
\end{proof}

At weight seven and depth four the imaginary balanced words have three
$z$'s, three $x$'s, and one $y$. There are $50$ such admissible words:
$140$ unrestricted permutations, minus $60$ beginning with $x$, minus
$60$ ending with $z$, plus $30$ satisfying both exclusions. Exactly
four are fixed by $C$. All fixed signs are positive. Their anti-invariant
space therefore has dimension $(50-4)/2=23$.
This completeness statement concerns the unmultiplied linear span.
The full ideal intersection has dimension $582$ by
\eqref{depthcay:eq:intersection}; products must be included to obtain
that larger object. It does not assert that a previously enumerated
collection of bounded products was already complete.

\subsection{Consequences for the frozen \texorpdfstring{$S_6$}{S6} search and verification}

For any specified family of additional exact relation rows
$\mathscr R\subset\mathcal A$, Theorem~\ref{depthcay:thm:projector}
gives the exact equivalence
\begin{equation}\label{depthcay:eq:combined}
 F\in\mathcal J+\operatorname{span}_{\mathbb Q}\mathscr R
 \quad\Longleftrightarrow\quad
 \Pi_{\uparrow}F\in
 \operatorname{span}_{\mathbb Q}\{\Pi_{\uparrow}R:R\in\mathscr R\}.
\end{equation}
If $F$ and the additional rows have depth at most $D$, both sides of
the projected linear-algebra problem stay in that same depth bound.
This is a concrete way to impose the complete Cayley ideal before
adding a prescribed double-shuffle or distribution family.

For a reproducible example, freeze the inherited word polynomial
\cite[Eq.~\texttt{cay:eq:target}]{depth:PriorProjection}
\begin{align*}
F={}&-179712000z^5aa+485683200z^5a\bar a
       -36864000z^3az^2a+401080320zaz^4a\\
 &-\frac{47600087040}{61}z^6a
       +11750400(za\shuffle z^4c)
       +109347840(z^3a\shuffle z^2c)\\
 &-971366400(z^5a\shuffle b).
\end{align*}
All powers inside words are concatenation powers. The imaginary part
of $H(F)$ is $485683200$ times the frozen conjectural $S_6$ residual.
The oriented normal form of $F^-=(F-KF)/2$ has $30$ nonzero
adapted-word coefficients, all
of depth at most two. The previous averaged output has $3444$ original-word
coordinates. The coordinate systems and representatives differ; these
support counts are descriptions of the supplied forms, not minimal-support
theorems. The new output is still nonzero. For example,
\[
 [z^5yh]\,\Pi_{\uparrow}(F^-)= -166348800,
 \qquad F^-=(F-KF)/2.
\]
Its full rational vector is supplied in
\texttt{s6\_oriented\_normal\_form.json}. This remains an obstruction
only for the Cayley ideal, already known without the depth refinement.
It neither proves nor disproves the numerical $S_6$ identity.

The accompanying code uses integers and rational fractions. It compares
the consecutive-cut Eulerian formula with its independent permutation
formula, verifies the endpoint retraction and all projector identities
on all $780$ words of weights one through four, and checks $2150$
shuffle products of total weight at most four. An independent
fraction-free integer elimination builds the entire ideal generated by
all Cayley rows and their products through weight four. At every depth
its ranks agree with \eqref{depthcay:eq:Hilbert}--\eqref{depthcay:eq:HilbertK}.
These finite audits check implementation and conventions. The proofs
above establish the statements in every weight.

A further useful computation is now well specified: apply the oriented
projector to the bounded weight-seven double-shuffle and distribution
rows, then search for the target in their rational span. The projected
problem can use the $1964$-dimensional image in
\eqref{depthcay:eq:weightseven}, without adding artificial higher-depth
coordinates merely to impose the Cayley ideal. An exact rational
certificate is still required before changing the conjecture's status.
